% Capitulo 8: listas con proposito.
\section{Listas}\label{sec:listas}

Tres entornos bastan para casi todo lo que hay que enumerar, y este capítulo los
usa a los tres en el orden en que conviene aprenderlos: primero lo que no
tiene orden, luego lo que sí lo tiene y por último lo que necesita nombre.
Cada lista tiene su oficio. La viñeta enumera, la numerada ordena y la de
descripción define; usar cada una para lo suyo es lo que separa un
documento legible de un amasijo de puntos.

Lo que texforge evita instalar cabe en una viñeta, porque el orden no
importa y la lista puede crecer sin reescribir nada:

\begin{itemize}
  \item ni TeX Live ni MiKTeX;
  \item ni Inkscape ni Node para los diagramas;
  \item ni \texttt{minted} ni \texttt{listings} para el código.
\end{itemize}

El ciclo de trabajo, en cambio, pide números, porque el orden sí
importa y cada paso presupone el anterior:

\begin{enumerate}
  \item escribir el fuente en texto plano;
  \item revisar con \texttt{texforge check};
  \item formatear con \texttt{texforge fmt};
  \item componer con \texttt{texforge build};
  \item inspeccionar con \texttt{texforge pdf}.
\end{enumerate}

Y los conceptos con nombre propio piden una lista de descripción, donde el
término manda y la definición sigue:

\begin{description}
  \item[Preámbulo] la cabecera del \texttt{.tex}: clase, paquetes,
  paleta, cabeceras y cajas. Se escribe una vez y gobierna todo el
  documento.
  \item[Entorno] un bloque con sintaxis propia, como \texttt{mermaid} o
  \texttt{code}: texforge lo detecta al compilar y lo dibuja.
  \item[Etiqueta] el nombre que \texttt{\textbackslash label} o
  \texttt{label=} da a una sección, figura, tabla, ecuación o
  listado para citarlo con \texttt{\textbackslash ref}.
\end{description}

\begin{notebox}
  La regla práctica: si reordenar los puntos cambia el sentido, numera;
  si cada punto necesita su término, describe; en lo demás, viñetas.
\end{notebox}
